% Part: first-order-logic
% Chapter: completeness
% Section: outline

\documentclass[../../../include/open-logic-section]{subfiles}

\begin{document}

\iftag{FOL}
      {\olfileid[tr]{fol}{com}{out}}
      {\olfileid[tr]{pl}{com}{out}}

\olsection{İspatın Ana Hatları}

Tamlık teoreminin ispatı biraz karmaşıktır ve ilk okumada yolunu
kaybetmek kolaydır. Bu nedenle önce ispatın ana hatlarını verelim. İlk
adım, ispata giden bir yol görmemizi sağlayan bir bakış açısı
değişikliğidir. Tamlık, ``$\Gamma \Entails !A$ olduğunda $\Gamma
\Proves !A$ olur'' biçiminde düşünüldüğünde bir fikir bulmak bile zor
olabilir: $\Gamma \Proves !A$ olduğunu göstermek için !!a{derivation}
bulmamız gerekir ve $\Gamma \Entails !A$ varsayımı bu konuda bize
hiç yardımcı olacakmış gibi görünmez. Bazı ispat sistemlerinde doğrudan
!!a{derivation} kurmak mümkündür; bizse biraz farklı bir yol izleyeceğiz.
Gerekli bakış açısı değişikliği şudur: tamlık, ``$\Gamma$ tutarlıysa
sağlanabilirdir'' biçiminde de ifade edilebilir. Belki~$\Gamma$'daki
bilgiyi, onun tutarlı olduğu varsayımıyla birlikte kullanarak
\iftag{FOL}{!!a{structure}}{!!a{valuation}} kurabiliriz; bu,
$\Gamma$'daki her \iftag{FOL}{!!{sentence}}{!!{formula}} ifadesini doğru
kılar. Sonuçta nasıl bir
\iftag{FOL}{!!{structure}}{!!{valuation}} aradığımızı biliyoruz:
$\Gamma$'nın betimlediği bir tane!

$\Gamma$ yalnızca \iftag{FOL}{atomik !!{sentence}s}{!!{propositional variable}s}
içeriyorsa ona bir model kurmak kolaydır.\iftag{FOL}{
  Atomik !!{sentence}s ifadelerinin hepsinin
  $\Atom{P}{a_1,\dots,a_n}$ biçiminde ve $a_i$ ifadelerinin
  !!{constant}s olduğunu varsayalım.}{} Yapmamız gereken tek şey
\iftag{FOL}{%
  !!a{domain}~$\Domain{M}$ ile~$P$ için öyle bir yorum bulmaktır ki
  $\Sat{M}{\Atom{P}{a_1,\ldots,a_n}}$ olsun. Bu hiç de zor değildir:
  $\Domain{M}=\Nat$, $\Assign{\Obj c_i}{M}=i$ olsun ve
  $\Atom{P}{a_1,\ldots,a_n}\in\Gamma$ olan her durumda $\tuple{k_1,
    \dots,k_n}$ n-lisini $\Assign{P}{M}$ içine koyalım; burada $k_i$,
  $a_i$ sabit simgesinin indisidir (yani $a_i \ident \Obj c_{k_i}$).}{%
  her $p\in\Gamma$ için $\pSat{v}{p}$ olacak bir
  !!a{valuation}~$\pAssign{v}$ bulmaktır. $\pAssign{v}(p)=\True$ ancak
  ve ancak $p\in\Gamma$ ise olsun.}

Şimdi $\Gamma$'nın atomik bir~$!B$ için $\lnot !B$ !!{formula}
ifadesini içerdiğini varsayalım. $\iftag{FOL}{\Struct{M}}{\pAssign{v}}$
kuruluşunun $\lnot !B$ ifadesini doğru kılmamızı engellemesinden
kaygılanabiliriz. Fakat burada~$\Gamma$'nın tutarlılığı devreye girer:
$\lnot !B\in\Gamma$ ise $!B\notin\Gamma$ olmalıdır; aksi halde
$\Gamma$ tutarsız olurdu. $!B\notin\Gamma$ olduğunda
$\iftag{FOL}{\Struct{M}}{\pAssign{v}}$ kuruluşumuza göre
$\iftag{FOL}{\Sat/{M}}{\pSat/{v}}{!B}$, dolayısıyla
$\iftag{FOL}{\Sat{M}}{\pSat{v}}{\lnot !B}$ olur. Buraya kadar sorun yok.

$\Gamma$ karmaşık, atomik olmayan formüller içeriyorsa ne olur?
Örneğin $!A\land !B$ içerdiğini varsayalım. Bunu doğru kılmak için hem
$!A$ hem de $!B$~$\Gamma$'daymış gibi ilerlemeliyiz. $!A\lor !B\in
\Gamma$ ise bunlardan en az birini doğru kılmamız, yani en az biri
$\Gamma$'daymış gibi ilerlememiz gerekir.

Bu gözlemler şu fikri akla getirir:~$\Gamma$'ya ek !!{formula}s
ekleyelim; bunlar (a)~ortaya çıkan kümeyi tutarlı tutacak,
(b)~olası her atomik !!{sentence}~$!A$ için ortaya çıkan
kümenin ya $!A$'yı ya da $\lnot !A$'yı içermesini sağlayacak ve
(c)~küme $!A\land !B$ içerdiğinde hem $!A$ hem $!B$'yi, $!A\lor !B$
içerdiğinde bunlardan en az birini, vb. içerecektir. Bunu (gerekirse
sonsuza dek) sürdürürüz.
Eklenen bütün !!{formula}s kümesine~$\Gamma^*$ diyelim. Yukarıdaki
kuruluşumuz bize öyle bir
\iftag{FOL}{!!a{structure}~$\Struct{M}$}{!!a{valuation}~$\pAssign{v}$}
verir ki bunun~$\Gamma^*$'daki bütün tümceleri sağladığını tümevarımla
ispatlayabiliriz. $\Gamma\subseteq\Gamma^*$ olduğundan, yapı ya da
değerleme~$\Gamma$'daki bütün tümceleri de sağlar. (a) ile~(b)'yi
güvence altına almanın yeterli olduğu ortaya çıkacaktır. (b) koşulunu
sağlayan bir tümceler kümesine \emph{tam} denir. Dolayısıyla görevimiz,
tutarlı~$\Gamma$ kümesini tutarlı ve tam bir~$\Gamma^*$ kümesine
genişletmektir.

\iftag{FOL}{%
Bu planda bir güçlük vardır: $\lexists[x][!A(x)]\in\Gamma$ ise bir
!!{constant}~$c$ seçip süreç içinde $!A(c)$ ekleyebilmeyi umarız. Fakat
bunu her zaman yapabileceğimizi nereden biliyoruz? Dilimizde yalnızca
birkaç !!{constant}s bulunabilir ve bunların her biri için $\lnot
!A(c)\in\Gamma$ olabilir. $!A(c)$ ifadesini de ekleyemeyiz; çünkü bu,
kümeyi tutarsız kılar ve~$\Struct{M}$ yapısının $!A(c)$ mi yoksa
$\lnot !A(c)$ mi doğru kılması gerektiğini bilemeyiz. Dahası, $\Gamma$
hiçbir sabit simgesi bulunmayan bir dildeki tümcelerden oluşabilir
(örneğin küme kuramının dilinde).

Bu sorunun çözümü, başlangıçta sonsuz sayıda sabit ve bunları
niceleyicilere doğru biçimde bağlayan tümceler eklemektir. (Elbette
bunun bir tutarsızlık doğuramayacağını doğrulamamız gerekir.)

İlk kuruluşumuz, atomik tümcelerde yalnızca !!{constant}s varsa iyi
çalışır. Ancak dilde !!{function}s da bulunabilir. Bu durumda her şeyin
çalışmasını sağlamak için bu !!{function}s simgelerine~$\Nat$ üzerinde
doğru fonksiyonları atamak zor olabilir. Bu nedenle başka bir hile
kullanırız: $\Obj c_i$ simgesini yorumlamak için $i$'yi kullanmak yerine
!!{constant}s kümesinin kendisini evren olarak alırız. Böylece
$\Struct M$ her !!{constant} simgesini kendisine atayabilir:
$\Assign{\Obj c_i}{M}=\Obj c_i$. Hatta neden sonuna kadar gitmeyelim?
$\Domain{M}$ dilin bütün \emph{terimleri} olsun. Bu durumda
!!{function}s simgelerine, argümanları ve değerleri terimler olan
fonksiyonlar atamanın apaçık bir yolu vardır:
!!{function}~$\Obj f^n_i$ simgesine,
girdi olarak verilen $n$ tane $t_1$, \dots,~$t_n$ teriminden
$\Obj f^n_i(t_1,\dots,t_n)$ terimini üreten fonksiyonu atarız.

Yapbozun son parçası~$\eq$ ile ne yapacağımızdır. !!{predicate}~$\eq$
sabit bir yoruma sahiptir: $\Sat{M}{\eq[t][t']}$ ancak ve ancak
$\Value{t}{M}=\Value{t'}{M}$ ise. Bir~$t$ teriminin !!{value} olarak
kendisini verecek biçimde kurarsak bu !!{structure}, $t$ ile $t'$ aynı
terim olmadıkça $\eq[t][t']$ biçimindeki \emph{hiçbir} tümceyi doğru
kılmaz. Bu elbette bir sorundur; çünkü !!{function}s içeren bir dildeki
hemen her ilginç kuramda, $t$ ile $t'$ aynı terim olmadığı halde
$\eq[t][t']$ biçiminde teoremler bulunur (örneğin aritmetik kuramlarında
$\eq[(\Obj 0+\Obj 0)][\Obj 0]$). Sorunu çözmek için~$\Struct M$
yapısının evrenini değiştiririz: $\Domain{M}$ içindeki nesneler olarak
terimler yerine terim kümeleri kullanırız; her küme,~$\Gamma$'daki
tümcelerin eşit olmasını gerektirdiği bütün terimleri içerir. Örneğin
$\Gamma$ bir aritmetik kuramıysa bu kümelerden biri $\Obj 0$,
$(\Obj 0+\Obj 0)$, $(\Obj 0\times\Obj 0)$ vb. terimleri içerir.
$\Obj 0$ simgesine bu kümeyi atayacağız ve bu kümenin, içindeki bütün
terimlerin---örneğin $(\Obj 0+\Obj 0)$ teriminin de---değeri olduğu
ortaya çıkacaktır. Böylece $\eq[(\Obj 0+\Obj 0)][\Obj 0]$ tümcesi bu
gözden geçirilmiş !!{structure} içinde doğru olacaktır.}{}

İzleyeceğimiz yol şudur. Önce !!{complete} tutarlı kümelerin
özelliklerini inceleriz. Özellikle !!a{complete} tutarlı kümenin
$!A\land !B$ ifadesini ancak ve ancak hem $!A$ hem de~$!B$'yi
içerdiğinde; $!A\lor !B$ ifadesini ancak ve ancak bunlardan en az birini
içerdiğinde içerdiğini vb. ispatlarız (\olref[ccs]{prop:ccs}).
\iftag{FOL}{Ardından ``doygun'' tümce kümelerini tanımlayıp inceleriz.
  Doygun bir küme, niceleyicili her !!{sentence} ifadesini onun
  örnekleriyle bağlayan koşullu ifadeleri içerir
  (\olref[hen]{defn:henkin-exp}). Her tutarlı~$\Gamma$ kümesinin her
  zaman doygun bir~$\Gamma'$ kümesine genişletilebildiğini gösteririz
  (\olref[hen]{lem:henkin}). Bir küme tutarlı, doygun ve !!{complete}
  ise ayrıca \iftag{prvEx}{$\lexists[x][!A(x)]$ ifadesini ancak ve ancak
    bir kapalı~$t$ terimi için $!A(t)$ ifadesini içerdiğinde}{}
  \iftag{defEx,defAll}{}{ ve }\iftag{prvAll}{$\lforall[x][!A(x)]$
    ifadesini ancak ve ancak bütün kapalı~$t$ terimleri için $!A(t)$
    ifadesini içerdiğinde}{} içerme özelliğine sahiptir
  (\olref[hen]{prop:saturated-instances}).}{}
Sonra \iftag{FOL}{doygun}{} tutarlı
$\iftag{FOL}{\Gamma'}{\Gamma}$ kümesini
\iftag{FOL}{doygun, tutarlı ve}{tutarlı ve} !!{complete} bir
$\Gamma^*$ kümesine genişletebileceğimizi gösteririz
(\olref[lin]{lem:lindenbaum}).
\iftag{FOL}{terim modelimiz~$\Struct M(\Gamma^*)$}{!!{valuation}
  ifademiz~$\pAssign v(\Gamma^*)$} bu~$\Gamma^*$ kümesinden
tanımlanacaktır. \iftag{FOL}{Terim modelinin evreni kapalı terimler
  kümesidir ve !!{predicate}s simgelerinin yorumu,~$\Gamma^*$ içindeki
  atomik !!{sentence}s tarafından}{Değerleme,~$\Gamma^*$ içindeki
  !!{propositional variable}s tarafından} belirlenir
(\olref[mod]{defn:termmodel}). \iftag{FOL}{Doygun,}{} tam ve tutarlı
kümelerin özelliklerini kullanarak
$\iftag{FOL}{\Sat{M(\Gamma^*)}}{\pSat{v(\Gamma^*)}}{!A}$ ancak ve ancak
$!A\in\Gamma^*$ olduğunu (\olref[mod]{lem:truth}) ve özellikle
$\iftag{FOL}{\Sat{M(\Gamma^*)}}{\pSat{v(\Gamma^*)}}{\Gamma}$ olduğunu
gösteririz. \iftag{FOL}{Son olarak~$\Gamma$ kümesi~$\eq$ de içeriyorsa
  terim modelinin nasıl tanımlanacağını ele alırız
  (\olref[ide]{defn:term-model-factor}) ve modelin~$\Gamma^*$ kümesini
  sağladığını gösteririz (\olref[ide]{lem:truth}).}{}

\end{document}
